% Abhakseite „Das kann ich“ – Zone nur im Gesamt (\mitzone)
%% TODO Ich-kann-Sätze: Platzhalter wie in den Aufgabendateien
\begin{abhakseite}
\ifdefined\mitzone
\abhakgruppe{Kennst du schon}
\abhak{1}{Punktprobe an Geraden- und Ebenengleichungen (einsetzen, Gleichung prüfen)}
\abhak{2}{Normalenvektor ablesen und die Paarregeln anwenden (parallel: Skalarprodukt null; senkrecht: kollinear)}
\abhak{3}{Kollinearität prüfen und ein Gleichungssystem auf Widerspruch untersuchen}
\abhak{4}{Lineare und quadratische Gleichungen sowie Ungleichungen im Parameter lösen (quadrieren, Beträge, Vorzeichenfälle)}
\abhak{5}{Winkelformeln (Kosinus für Vektorpaare und Ebenen, Sinus für Gerade–Ebene) und die Hessesche Normalform anwenden}
\abhak{6}{Schnittpunkte einer Ebene mit Kanten und Achsen berechnen und Schrägbilder lesen}
\abhak{7}{Normalenvektor ablesen und die Paarregeln anwenden (parallel: Skalarprodukt null; senkrecht: kollinear) – Fehler finden}
\abhak{8}{Normalenvektor ablesen und die Paarregeln anwenden (parallel: Skalarprodukt null; senkrecht: kollinear) – selbst rechnen}
\fi
\abhakgruppe{Einheit 1 · Die Schar als Familie}
\abhak{9}{Was wählt der Parameter?}
\abhak{10}{Die Schar als Familie}
\abhak{11}{Die Schar als Familie – Fehler finden, Begründen}
\abhakgruppe{Einheit 2 · Parameter aus Lagebedingungen}
\abhak{12}{Parameter aus Lagebedingungen}
\abhak{13}{Parameter aus Lagebedingungen – Fehler finden, Begründen}
\abhakgruppe{Einheit 3 · Parameter aus Maßbedingungen}
\abhak{14}{Parameter aus Maßbedingungen}
\abhak{15}{Parameter aus Maßbedingungen – Fehler finden, Begründen}
\abhakgruppe{Einheit 4 · Scharen am Körper}
\abhak{16}{Scharen am Körper}
\abhak{17}{Scharen am Körper – Fehler finden, Begründen}
\end{abhakseite}
